\section{IA de nivel empresarial: los token públicos llevan a una competencia feroz; los datos privados tienen valor}

Las dos secciones anteriores explicaron por qué Minglüe (明略) cuenta con puntos de acceso a datos, experiencia de entrega y resistencia operativa; esta sección entra en el núcleo técnico de todo el episodio, y la pregunta pasa a ser: cómo se convierte esa acumulación, en la era de la IA, en un Agentic Model de nivel empresarial. Wu Minghui (吴明辉) considera que las empresas que entrenan modelos base con datos públicos y cuyo modelo de negocio es vender token enfrentarán una competencia feroz, que fácilmente terminará reduciéndose al costo de la electricidad. Esto se debe a que los datos públicos y las capacidades de los modelos generales tienden a converger, y el precio del token bajará. La diferenciación de la IA empresarial no está en volver a entrenar un modelo público, sino en los datos privados, los procesos de la industria y los sistemas de herramientas reales.

\begin{figure}[H]
\centering
\includegraphics[width=0.94\textwidth]{private-data-value.png}
\caption{Valor diferencial de los datos privados: vender token con datos públicos lleva a una competencia feroz; los datos privados pueden generar valor para la industria. Diagrama conceptual propio, elaborado a partir de la conversación de 01:45:01--02:06:40.}
\end{figure}

\begin{knowledgebox}{Lectura de la figura: por qué los datos privados generan un sobreprecio}
Las capacidades entrenadas con datos públicos se vuelven cada vez más una mercancía; los datos privados de una empresa incluyen registros de clientes, procesos, transacciones, contratos, conversaciones, equipos y riesgos. Solo cuando el modelo entra en ese contexto puede generar resultados diferenciados por los que la empresa esté dispuesta a pagar.
\end{knowledgebox}

\begin{warningbox}{Los «datos privados» no consisten simplemente en subir los archivos de la empresa}
Lo verdaderamente difícil de los datos privados de una empresa está en los permisos, la actualización en tiempo real, la estructuración, los datos sucios, el significado de negocio y los límites de responsabilidad. Meter los documentos en una base de datos vectorial es solo el primer paso; lograr que el Agent entienda correctamente quién puede ver, quién puede modificar y qué resultado es entregable es el verdadero umbral de la IA empresarial.
\end{warningbox}

\begin{knowledgebox}{Arquitectura de cuatro capas de la IA empresarial}
\begin{tabularx}{\textwidth}{p{0.20\textwidth}p{0.32\textwidth}X}
\toprule
Capa & Función & Forma típica de falla \\
\midrule
Capa de datos & Integrar datos privados, permisos e historial & Datos sucios, permisos erróneos, contexto incompleto. \\
Capa de modelo & Comprender, razonar y generar acciones candidatas & Solo sabe resumir, no entiende las restricciones del negocio. \\
Capa de herramientas & Invocar sistemas de negocio, bases de datos e interfaces de procesos & No puede escribir de vuelta en el sistema; se queda en sugerencias. \\
Capa de gobernanza & Auditoría, revisión, responsabilidad y retroalimentación & La empresa no se atreve a delegar y no es posible desplegar a escala. \\
\bottomrule
\end{tabularx}
\end{knowledgebox}

\begin{knowledgebox}{Cadena de conversión del modelo general a la entrega empresarial}
\begin{tabularx}{\textwidth}{p{0.22\textwidth}p{0.32\textwidth}X}
\toprule
Eslabón de conversión & Pregunta que debe responder & Valor de la experiencia tipo Minglüe \\
\midrule
Semántica de la industria & ¿A qué se refiere exactamente la empresa cuando dice «cliente», «riesgo» o «prospecto»? & La larga trayectoria de entrega ToB acumuló semántica de la industria y comprensión del entorno real del cliente. \\
Estado del proceso & ¿La tarea está en preventa, cumplimiento del contrato, cobranza o posventa? & Los datos de CRM, ventas y conversaciones de atención al cliente aportan la trayectoria del proceso. \\
Medición de resultados & ¿La acción mejora de verdad los ingresos, los costos o el riesgo? & La experiencia de Miaozhen (秒针) y en análisis de datos pone énfasis en que todo sea medible y auditable. \\
Adopción organizacional & ¿Quién aprueba, quién revisa, quién asume las consecuencias? & La operación ToB formó la conciencia de permisos, responsabilidad y entrega. \\
\bottomrule
\end{tabularx}
\end{knowledgebox}

Aquí también hay implícito un juicio sobre el modelo de negocio: si una empresa de modelos solo vende token públicos, el precio se verá presionado continuamente a la baja por el aumento de la oferta y la optimización de la inferencia; si una empresa de servicios empresariales convierte los token en acciones de negocio con precio asignable, el precio unitario ya no depende solo del costo de cómputo. Por ejemplo, con el mismo consumo de un millón de token, generar un resumen ordinario y descubrir un riesgo contractual grave tienen un valor comercial completamente distinto. La unidad de precio de la IA empresarial podría terminar pasando del token a las tareas, los riesgos, el aumento de ingresos y el ahorro de costos.

\begin{importantbox}{Del precio por token al precio por resultado}
La unidad de cobro de las API de modelos públicos suele ser el token; la unidad de cobro ideal de la IA empresarial debería ser el resultado. Solo cuando la salida del modelo puede entrar en el workflow y producir beneficios de negocio verificables, la empresa pagará por «resultados» y no por «número de palabras».
\end{importantbox}

\subsection{El juego numérico del mundo real}

El apartado anterior trató los datos privados; este explica por qué es difícil convertir directamente los datos empresariales en capacidades del modelo. Wu Minghui llama a los servicios empresariales el juego numérico del mundo real. El go se presta muy bien a la IA porque el tablero es transparente, las reglas son completas y el context es total; el trabajo empresarial es lo contrario: el contexto está incompleto, la información está dispersa, los objetivos son ambiguos, los procesos cruzan sistemas y los resultados deben responder ante clientes reales y ante las finanzas. La dificultad de la IA empresarial no es «si sabe responder», sino si puede convertir un contexto incompleto en acciones ejecutables.

\begin{importantbox}{Fórmula de la dificultad de la IA empresarial}
\[
\text{Enterprise AI Value} \approx f(\text{Private Data}, \text{Workflow}, \text{Tool Use}, \text{Trust})
\]
Donde \(\text{Private Data}\) es el contexto privado de la empresa, \(\text{Workflow}\) son los procesos de negocio, \(\text{Tool Use}\) es la capacidad de invocar sistemas y herramientas, y \(\text{Trust}\) es la confianza de que todo sea auditable, controlable y entregable.
\end{importantbox}

\begin{knowledgebox}{Términos clave: las cuatro restricciones de la IA empresarial}
\begin{tabularx}{\textwidth}{p{0.20\textwidth}p{0.32\textwidth}X}
\toprule
Restricción & Significado & Consecuencia para el producto \\
\midrule
Private Data & Datos de clientes, procesos, transacciones y conversaciones exclusivos de la empresa & Determina si el modelo entiende el entorno real del negocio. \\
Workflow & Procesos reales que cruzan departamentos y sistemas & Determina si el Agent puede pasar de la sugerencia a la entrega. \\
Tool Use & Invocar herramientas como CRM, ERP, sistemas de tickets y bases de datos & Determina si la IA puede cambiar el estado del negocio. \\
Trust & Permisos, auditoría, revisión y límites de responsabilidad & Determina si la empresa se atreve a confiarle tareas a la IA. \\
\bottomrule
\end{tabularx}
\end{knowledgebox}

Por eso «el juego numérico del mundo real» es más difícil que el go. El espacio de estados del go es enorme, pero el tablero está completo, las reglas son claras y el resultado se puede determinar; el espacio de estados del negocio empresarial no solo es grande, sino que muchas variables están ocultas en la mente de las personas, en las relaciones con los clientes, en los contratos históricos y en los procesos entre sistemas. Para que un Agent pueda trabajar aquí, primero debe convertir el conocimiento tácito en memory aprovechable por el sistema y luego combinar esa memory con la invocación de herramientas.

\subsection{Agentic Model: de las capacidades del modelo a la ejecución de flujos de trabajo}

Este apartado explica el Agentic Model. En la entrevista, Wu Minghui ve al Agent o al Tool Use como el nuevo vínculo en el ámbito de los servicios empresariales. El modelo ya no solo responde preguntas, sino que se conecta a los datos, comprende la tarea, invoca herramientas, completa el proceso y devuelve resultados auditables. Esto se relaciona con el Agent loop de EP139 y el Agent OS de EP118, pero el escenario empresarial pone más énfasis en la implementación privada, los permisos, la auditoría y los límites de responsabilidad.

\begin{figure}[H]
\centering
\includegraphics[width=0.96\textwidth]{enterprise-agentic-model.png}
\caption{Agentic Model de nivel empresarial: de los datos privados a la invocación de herramientas y, después, a la ejecución de los flujos de trabajo de la empresa. Diagrama conceptual propio, elaborado a partir de la conversación de 01:45:01--03:00:00.}
\end{figure}

\begin{knowledgebox}{Lectura de la figura: un Agent empresarial no es un chatbot}
De izquierda a derecha, el Agent empresarial primero debe obtener los datos privados y el contexto de la empresa, luego usar el modelo para comprender y razonar, después invocar mediante Tool Use sistemas como CRM, ERP, sistemas de tickets y bases de datos, y por último entrar en un flujo de trabajo auditable. Si falta cualquier eslabón, es muy difícil que se convierta en un producto de nivel empresarial.
\end{knowledgebox}

\subsection{DeepMiner: extraer de los datos empresariales pistas para la acción}

El Agentic Model resuelve el «cómo ejecutar»; este apartado usa DeepMiner para explicar «dónde se descubren las tareas». DeepMiner aparece poco en la entrevista, pero es una pista para entender la dirección de los nuevos productos de Minglüe. Representa una forma de producto de IA empresarial: no se limita a preguntas y respuestas, sino que integra datos de múltiples sistemas, extrae reglas, riesgos y oportunidades, genera hallazgos y luego ejecuta las acciones siguientes mediante un Agent o herramientas.

\begin{figure}[H]
\centering
\includegraphics[width=0.96\textwidth]{deepminer-product-loop.png}
\caption{Ciclo cerrado del producto DeepMiner: usar la IA para extraer reglas, riesgos y pistas para la acción de los datos de la empresa. Diagrama conceptual propio, elaborado a partir de la conversación de 02:31:58--03:36:59.}
\end{figure}

\begin{knowledgebox}{Lectura de la figura: de la minería de datos a la entrega Agentic}
El BI y el análisis de datos tradicionales se orientan más a mostrar hallazgos; un producto tipo DeepMiner, si se conecta a un Agent, puede ir más allá de «descubrir el problema» y llegar a «proponer acciones», «invocar sistemas» y «registrar resultados». Este es precisamente el punto decisivo para que la IA empresarial pase de ser una herramienta de análisis a ser un sistema de producción.
\end{knowledgebox}

\begin{importantbox}{El cambio decisivo de los productos tipo DeepMiner}
El análisis de datos de antes solía quedarse en dashboard e informes; un producto Agentic de nivel empresarial debe responder además: qué anomalía se descubrió, qué acción debe tomarse, quién aprueba la acción, qué sistema se invoca y cómo regresa el resultado.
\end{importantbox}

\begin{knowledgebox}{De Copilot a Agentic System}
\begin{tabularx}{\textwidth}{p{0.22\textwidth}p{0.34\textwidth}X}
\toprule
Etapa & Forma típica & Requisitos del lado de la empresa \\
\midrule
Copilot & La persona pregunta, la IA responde y la persona copia el resultado & Asistencia de bajo riesgo y alta frecuencia; mejora sobre todo la eficiencia individual. \\
Workflow Agent & La IA lee datos, llena formularios e inicia procesos & Debe integrar permisos, herramientas y aprobaciones. \\
Agentic System & Varios Agent colaboran para completar tareas & Requiere registros, orquestación, recuperación ante fallas y responsabilidad sobre los resultados. \\
Trusted Agentic Model & Asume de forma confiable tareas críticas de la empresa & Debe ser auditable, revisable y desplegable de forma privada o ejecutarse de forma controlada. \\
\bottomrule
\end{tabularx}
\end{knowledgebox}

\subsection{Resumen de la sección}

El núcleo de la IA empresarial está en los datos privados y los flujos de trabajo. Los token públicos se abaratarán; lo que realmente tiene un sobreprecio es el Agentic Model capaz de convertir el contexto de la empresa en resultados.
